\begin{helper}[Copying augmentations through the first missing-color vertex]
Let $\varphi$ be a proper partial edge-coloring of a finite loopless
multigraph, $\alpha\ne\beta$, and let $\psi$ be the interchange of
these colors on an arbitrary full two-color component rooted at a node
$z$. Suppose $R$ is obtained from the tree-sequence $S$ by exact
$\varphi$-Tashkinov augmentations. Let $X$ be elementary under
$\varphi$, with $V(R)\subseteq X$. Suppose $u\in X\setminus V(R)$
misses $\alpha$ and $v\in V(R)\setminus V(S)$ misses $\beta$.
Then there is an exact ordered initial segment $B$ of $R$, whose last
vertex is $v$, obtained from the same $S$ by $\psi$-Tashkinov
augmentations.
\end{helper}
\begin{proof}
By \noderef{TashkinovAugmentationSequence}, choose the finite chain
of exact augmentations from $S$ to $R$. Each one-step augmentation
adds one vertex at the end and preserves every previous vertex and
edge entry, by \noderef{TashkinovAugmentation} and
\noderef{TreeSequenceInitialSegment}. Every intermediate tree is an
exact initial segment of $R$ by
\noderef{TashkinovAugmentationSequenceInitialSegment}; its vertex set
is therefore contained in $V(R)$ by
\noderef{TreeSequenceInitialSegmentVertexSetSubset}.
Since $v$ is absent from $S$ but present in $R$, there is a first step
adjoining $v$. Write $B$ for its output. It is an exact initial segment
of $R$ and its last vertex is $v$: the old vertices are unchanged, so
the first appearance of $v$ must be the unique new vertex.

Consider any step $A\longrightarrow D$ up to and including this
step. In particular $v\notin V(A)$. The step certificate in
\noderef{TashkinovAugmentation} supplies the exact new edge identity
$f$, its old color $c$, and $c\in\overline\varphi(A)$.
By \noderef{TreeMissingColors}, choose $w\in V(A)$ with
$c\in\overline\varphi(w)$. We have $w\in V(R)\subseteq X$.
Since $u\notin V(R)$, $w\ne u$. Since $v\notin V(A)$, $w\ne v$,
including at the step adjoining $v$. Also $u,v\in X$.
Elementarity, in the precise pairwise-disjoint sense of
\noderef{ColoredVertexSetElementary}, now gives both exclusions:
if $c=\alpha$, the missing sets at $w,u$ would both contain
$\alpha$; if $c=\beta$, the missing sets at $w,v$ would both contain
$\beta$. Thus $c\notin\{\alpha,\beta\}$.

The definition \noderef{TwoColorComponentSwap} consequently fixes
the exact value of $f$, whether or not its edge-node lies in the
swapped component. In particular $\psi(f)=\varphi(f)=c$.
The unaffected-color clause of
\noderef{TwoColorComponentSwapMissingBehavior} says that $c$ remains
missing at the same $w$ under $\psi$. Thus
\noderef{TreeMissingColors} gives $c\in\overline\psi(A)$.
The length and exact-entry clauses are independent of the coloring,
so all clauses of \noderef{TashkinovAugmentation} hold for the same
step $A\longrightarrow D$ under $\psi$.

Apply this argument to every step through the one adjoining $v$.
Concatenating these steps with
\noderef{TashkinovAugmentationSequence} gives the required route
from the unchanged seed $S$ to the exact prefix $B$. No edge is replaced
by another edge with the same endpoints, and no preservation of the
whole missing set at an internal component vertex is asserted or used.
\end{proof}
